\begin{abstract}
Metric depth models must infer the camera from pixels to output metres. We use a digital zoom,
a crop followed by a resize, as a free and exact test of whether they do: the focal length is
multiplied by the zoom factor $s$ and the depth of every surviving pixel is unchanged. Two
exponents summarise a model's response: its depth scales like $s^{-\alpha}$ (ideal $\alpha=0$) and
its predicted focal like $s^{\beta}$ (ideal $\beta=1$). Neither needs ground truth. Across
\nImages{} images from KITTI, NYUv2 and DIODE, the widely used Depth Anything V2 has
$\alpha=\DavTwoSAlpha$, close to reading depth from apparent size alone. Models that predict their
own camera close only part of the gap ($\alpha\approx\DepthproAlpha$--\MogeTwoSAlpha) and under-scale the
focal ($\beta=\UnidepthTwoSBeta$--\DepthproBeta); Depth Pro lies on the line $\alpha+\beta=1$ that
apparent-size reasoning predicts (median \DepthproAlphaPlusBeta{} per image). We explain this
failure mode with a zoom--squash ambiguity: a zoomed photo is pixel-for-pixel the photo of the same
scene compressed along the viewing axis, which size cues cannot detect. Controls show the effect is
not caused by interpolation or detail loss. Finally, a self-consistency fine-tuning on unlabeled
images reduces $\alpha$ to \DavTwoSEftbKittiAlpha{} on KITTI and \DavTwoSEftbNyuAlpha{} on NYUv2,
once a blur-invariance term removes a ``blurry means zoomed'' shortcut that the plain objective
learns; we report its cost on unzoomed and unseen images.
\end{abstract}
